\section*{ВВЕДЕНИЕ}
\addcontentsline{toc}{section}{ВВЕДЕНИЕ}

Булевые функции играют ключевую роль в симметричной криптографии,
используясь как функции замены в блочных шифрах и компоненты генераторов
ключевого потока. Криптостойкость таких систем напрямую зависит от свойств
булевых функций: сбалансированности, нелинейности, дифференциального
поведения и алгебраической иммунности. Многие из этих требований
противоречат друг другу, что существенно усложняет конструирование функций
с оптимальным набором свойств.

\textbf{Актуальность темы} определяется тем, что центральным нелинейным
элементом современных блочных шифров является блок подстановки (S-блок),
стойкость которого полностью определяется свойствами составляющих его
булевых функций. Несмотря на то, что S-блоки действующих стандартов
удовлетворяют установленным криптографическим требованиям, ряд
теоретических вопросов остаётся открытым. Разработка инструмента, совмещающего анализ булевых свойств и направленный поиск подстановок с заданными характеристиками, создаёт
основу для воспроизводимых экспериментальных исследований в данной области.

\textbf{Целью работы} является разработка программного инструмента для
анализа криптографических свойств булевых функций и генерации S-блоков, удовлетворяющих заданным критериям стойкости
к дифференциальному и линейному криптоанализу.

Для достижения цели поставлены следующие \textbf{задачи}:
\begin{itemize}
    \item \textbf{изучить} криптографические характеристики булевых функций,
        используемые для оценки стойкости S-блоков;;
    \item \textbf{разработать} модуль вычисления булевых характеристик
        произвольного S-блока;
    \item \textbf{оценить} работоспособность инструмента на S-блоках
        известных шифров и оценить качество генерируемых подстановок.
\end{itemize}

\section{Современные блочные шифры и роль S-блоков}
\subsection{Общая структура блочных шифров}
\subsection{Обзор современных блочных шифров}
\indent Рассмотрим применение S-блоков в шести широко
применяемых блочных шифрах: \textit{AES}, ГОСТ~Р~34.12-2015 («Кузнечик»), \textit{Serpent},
\textit{Twofish}, \textit{SM4} и \textit{Camellia}.

\subsubsection{\textit{AES}}

В 1997 году Национальный институт стандартов и технологий США
инициировал открытый конкурс на замену устаревшего стандарта DES.
По итогам многолетней публичной экспертизы в 2001 году алгоритм Rijndael,
разработанный бельгийскими криптографами Йоаном Дайменом
и Винсентом Рейменом, был принят в качестве федерального
стандарта шифрования AES~\cite{fips197}.

AES является блочным шифром с размером блока 128~бит и поддержкой ключей
длиной 128, 192 и 256~бит. Структура алгоритма представляет собой
подстановочно-перестановочную сеть (\textit{SPN}). Каждый из 10, 12 или 14 раундов
(в зависимости от длины ключа) включает четыре преобразования:
\textit{SubBytes}, \textit{ShiftRows}, \textit{MixColumns} и \textit{AddRoundKey}.

AES использует единственный статичный 8-битный S-блок, представляемый
матрицей $16 \times 16$. Имеет нелинейность~112 при теоретическом максимуме~120, дифференциальную
равномерность~4 и алгебраическую степень~7.

% --- 2.2 Кузнечик ---
\subsubsection{ГОСТ~Р~34.12-2015}

Российский национальный стандарт блочного шифрования ГОСТ~Р~34.12-2015
был принят в 2015 году. Алгоритм «Кузнечик» разработан ФСБ России совместно
с компанией «ИнфоТеКС» и пришёл на смену ГОСТ~28147-89 («Магма»), чьи
S-блоки изначально имели засекреченное происхождение.~\cite{gost2015}.

«Кузнечик» оперирует блоками размером 128~бит при длине ключа 256~бит
и включает 10 раундов. Структура представляет собой подстановочно-перестановочную сеть.
Стандарт определяет один статичный 8-битный S-блок~$\pi$, заданный в виде
константной таблицы подстановки без раскрытия метода его построения.
Данный подход стал предметом научной дискуссии: в работе~\cite{biryukov2016}
было показано, что S-блок «Кузнечика» может быть представлен в виде
композиции определённых степеней в $\mathrm{GF}(2^8)$, что, по мнению ряда
исследователей, потенциально указывает на наличие скрытой структуры.
Нелинейность S-блока составляет~104, дифференциальная равномерность 8.
По обоим показателям он уступает S-блоку \textit{AES}, однако свойства линейного слоя
стандарта компенсируют данное различие на уровне всего шифра.

\subsubsection{Serpent}

\textit{Serpent} был представлен в 1998 году как один из пяти финалистов конкурса
AES. Авторами алгоритма являются Росс Андерсон (Великобритания), Эли Бихам
(Израиль) и Ларс Кнудсен (Дания). По итогам голосования \textit{Serpent} занял второе
место, набрав наибольшее число голосов среди участников, ставивших приоритетом
максимальную безопасность в ущерб производительности~\cite{serpent1998}.

Serpent использует блок размером 128~бит и ключи длиной 128, 192 или 256~бит.
Алгоритм включает 32 раунда, что вдвое превышает количество раундов \textit{AES}.
Авторы сознательно придерживались консервативного подхода, обеспечивающего
значительный запас стойкости к будущим атакам. Эффективность программных
реализаций достигается за счёт техники «битовых срезов» (\textit{bitslicing}),
позволяющей обрабатывать несколько блоков параллельно.

Serpent использует набор из восьми различных 4-битных S-блоков, применяемых
циклически по раундам. Каждый из них реализует биективное отображение
$\{0,1\}^4 \to \{0,1\}^4$. S-блоки выбирались из перестановок, производных
от S-блоков DES, с оптимизацией по критериям дифференциальной и линейной
стойкости. Нелинейность каждого S-блока равна~4~-максимально достижимой
для 4-битных булевых функций. Благодаря 32 раундам совокупная стойкость
\textit{Serpent} к дифференциальному и линейному криптоанализу превосходит
соответствующие показатели AES.

\subsubsection{Twofish}

\textit{Twofish} был разработан в 1998 году командой под руководством Брюса Шнайера
(Counterpane Internet Security) при участии Джона Келси, Дага Вина,
Дэвида Вагнера и Криса Холла. Алгоритм вошёл в пятёрку финалистов конкурса
\textit{AES}~\cite{twofish1998}.

\textit{Twofish} оперирует блоками размером 128~бит при ключах 128, 192 и 256~бит,
выполняя 16 раундов в структуре сети Фейстеля. Отличительной особенностью
алгоритма является использование ключезависимых S-блоков: конкретные таблицы
подстановки формируются из ключевого материала на этапе \textit{KeySchedule}.
Кроме того, \textit{Twofish} применяет псевдо-адамарово преобразование (PHT)
и \textit{MDS}-матрицу в функции Фейстеля.

Шифр использует четыре 8-битных S-блока, каждый из которых строится как
композиция двух фиксированных 4-битных перестановок ($q_0$ и $q_1$)
с промежуточным наложением ключевого материала через XOR и умножением на
\textit{MDS}-матрицу. Таким образом, алгебраическая структура S-блоков является
фиксированной, тогда как конкретные значения таблиц зависят от ключа.
Такой подход представляет собой компромисс между полностью статичными
и полностью случайными ключезависимыми подстановками и теоретически затрудняет
атаки, предполагающие предварительное знание S-блоков.

\subsubsection{SM4}

\textit{SM4} — китайский государственный стандарт блочного шифрования. Алгоритм
был разработан Государственным бюро криптографии КНР (\textit{OSCCA}) и впервые
опубликован в 2006 году в контексте стандарта беспроводных локальных
сетей \textit{WLAN}. В 2016 году \textit{SM4} был принят в качестве национального
стандарта и активно продвигается как альтернатива \textit{AES}
в государственных и коммерческих системах Китая~\cite{sm4_2006}.

\textit{SM4} использует блок размером 128~бит при длине ключа 128~бит и выполняет
32 раунда в обобщённой структуре сети Фейстеля. Каждый раунд применяет
функцию~$T$, состоящую из нелинейного преобразования~$\tau$ (побайтовое
применение S-блока) и линейного преобразования~$L$. Алгоритм расширения
ключа построен по принципу, аналогичному основному шифру.


\textit{SM4} использует один статичный 8-битный S-блок, заданный таблицей подстановки.
Происхождение S-блока долгое время оставалось официально нераскрытым,
что порождало вопросы, схожие с дискуссией вокруг «Кузнечика». Последующий
криптографический анализ показал, что S-блок \textit{SM4} построен по принципу,
аналогичному \textit{AES}: вычисление мультипликативной инверсии в $\mathrm{GF}(2^8)$
с использованием иного неприводимого полинома в сочетании с аффинным
преобразованием. Нелинейность S-блока составляет~112, дифференциальная
равномерность~4, что идентично параметрам \textit{AES}.


\subsubsection{Camellia}

\textit{Camellia }разработан совместно японскими компаниями \textit{NTT} и \textit{Mitsubishi Electric}
в 2000 году. Алгоритм прошёл европейский конкурс \textit{NESSIE}
и японский конкурс \textit{CRYPTREC}, получив статус рекомендованного шифра.
\textit{Camellia} включён в стандарты \textit{TLS}, \textit{IPsec} и широко применяется
в японской государственной и коммерческой инфраструктуре~\cite{camellia2001}.

\textit{Camellia }использует блок размером 128~бит при ключах 128, 192 или 256~бит.
Структура~ — сеть Фейстеля с 18 раундами (для ключа 128~бит) или 24 раундами
(192/256~бит). Характерной особенностью является наличие
$\mathrm{FL}/\mathrm{FL}^{-1}$-функций, вставляемых каждые 6 раундов:
они образуют дополнительный нелинейный слой, существенно усложняющий полный
дифференциальный и линейный криптоанализ. Алгоритм изначально проектировался
с учётом эффективности как программных, так и аппаратных реализаций.


\textit{Camellia} использует четыре 8-битных S-блока: один базовый ($s_1$) и три
производных ($s_2$, $s_3$, $s_4$), полученных циклическим сдвигом выходного
байта базового S-блока на 1~бит. Базовый S-блок построен аналогично \textit{AES}:
инверсия в $\mathrm{GF}(2^8)$ с иными параметрами поля и аффинным
преобразованием с другой матрицей. Нелинейность составляет~112,
дифференциальная равномерность~4. Использование четырёх родственных
S-блоков вместо одного повышает диффузию и усложняет алгебраический анализ,
сохраняя при этом компактность аппаратной реализации.

\subsection{Сравнительный анализ S-блоков современных стандартов}
Рассмотренные блочные шифры демонстрируют различные подходы к построению
нелинейного слоя. В \textit{AES}, \textit{SM4} и «Кузнечике» используется
единственный статичный 8-битный S-блок. В \textit{Camellia} применяются
четыре родственные 8-битные подстановки, полученные из базового S-блока.
Алгоритм \textit{Serpent} использует набор из восьми различных 4-битных
S-блоков, циклически применяемых по раундам. В \textit{Twofish} реализован
ключезависимый подход, при котором конкретные таблицы подстановки
формируются из ключевого материала.

С точки зрения количественных криптографических характеристик
наблюдаются устойчивые закономерности. Для 8-битных S-блоков
современных стандартов характерна нелинейность 112
(\textit{AES}, \textit{SM4}, \textit{Camellia}), что близко к
теоретическому максимуму 120. В «Кузнечике» нелинейность составляет 104.
Для 4-битных S-блоков \textit{Serpent} нелинейность равна 4,
что является максимально возможным значением для данной размерности.

Дифференциальная равномерность 8-битных подстановок в \textit{AES},
\textit{SM4} и \textit{Camellia} равна 4, что является оптимальным
значением для биективных отображений $\{0,1\}^8 \to \{0,1\}^8$.
В «Кузнечике» данный показатель равен 8. Для 4-битных S-блоков
\textit{Serpent} максимальная дифференциальная вероятность также
соответствует оптимальным значениям для выбранной размерности.

Алгебраическая степень 8-битных S-блоков рассматриваемых стандартов,
как правило, равна 7, что близко к максимально достижимому значению
для биективных отображений данной размерности.

Таким образом, анализ действующих стандартов позволяет выделить
практические ориентиры для 8-битных S-блоков современного уровня
криптографической стойкости: нелинейность не ниже 104,
дифференциальная равномерность не выше 8 и алгебраическая степень
не ниже 7. Полученные значения демонстрируют практический уровень
требований к S-блокам современных стандартов.

\subsection{Постановка задачи исследования и обоснование разработки инструмента}
Проектирование нелинейных компонентов блочных шифров S-блоков
представляет собой самостоятельную криптографическую задачу, актуальность
которой определяется несколькими факторами. 

Для биективных булевых функций размерности $8 \times 8$ точные границы достижимых
значений нелинейности при ограниченной дифференциальной равномерности
до сих пор не установлены, что делает экспериментальный поиск в данном
пространстве теоретически значимым. 

S-блоки действующих стандартов (\textit{AES}, \textit{SM4}, \textit{Camellia}) построены на основе инверсии в конечном
поле $\mathrm{GF}(2^8)$ и обладают компактным алгебраическим описанием,
что потенциально упрощает алгебраический анализ; поиск подстановок без
явной алгебраической структуры при сопоставимых метриках представляет
самостоятельный интерес.

Дополнительную сложность представляет наличие теоретических
противоречий между криптографическими характеристиками булевых функций: оптимизация одной характеристики нередко приводит к ухудшению другой. В результате проектирование S-блока
представляет собой задачу многокритериальной оптимизации, в которой
требуется найти допустимый компромисс между совокупностью количественных
показателей.

Несмотря на наличие устоявшихся стандартов, отсутствует общедоступный
универсальный инструмент, позволяющий:
\begin{itemize}
    \item вычислять полный набор булевых характеристик произвольного S-блока;
    \item анализировать их устойчивость к криптоанализу;
    \item осуществлять поиск S-блоков с заданными пороговыми значениями.
\end{itemize}

Отсутствие такого инструментария затрудняет воспроизводимость
исследований, сопоставимость результатов и экспериментальное
проектирование новых подстановок.

В связи с изложенным целью настоящей работы является разработка
программного инструмента, совмещающего функции анализатора булевых
свойств S-блоков и генератора подстановок, удовлетворяющих заданным
криптографическим критериям. Пороговые значения характеристик
определяются с учётом практических ориентиров, выявленных при анализе
действующих стандартов.
\subsection*{ВЫВОДЫ}
В первой главе рассмотрены современные блочные шифры и проанализированы
используемые в них S-блоки. Проведён сравнительный анализ количественных характеристик S-блоков,
выделены практические ориентиры,характерные для стандартов современного уровня стойкости.
\addcontentsline{toc}{subsection}{ВЫВОДЫ}


\newpage
\section{Стойкость к линейному криптоанализу}
\begin{definition}
S-блок (substitution box, блок подстановки) — это функция, выполняющая нелинейное преобразование в блочных шифрах.
S-блок размера $n\times m$ — это отображение:
\[S : \mathbb{F}_2^n \to \mathbb{F}_2^m\]
\end{definition}
\subsection{Линейный криптоанализ}

Суть линейного криптоанализа заключается в линейной аппроксимации единственного нелинейного компонента блочного шифра — S-блока. Для определения корреляции входных и выходных битов рассматривают следующее выражение:
\[
    \alpha_1x_1 \oplus \alpha_2x_2 \oplus \dots \oplus \alpha_nx_n \oplus \beta_1y_1 \oplus \beta_2y_2 \oplus \dots \oplus \beta_my_m = 0,
\]
где $\alpha = (\alpha_1, \alpha_2, \dots, \alpha_n) \in \mathbb{F}_2^n$ — маска входных битов; \\
\indent$\beta = (\beta_1, \beta_2, \dots, \beta_m) \in \mathbb{F}_2^m$ — маска выходных битов; \\
\indent$x = (x_1, \dots, x_n)$ — входной вектор; \\
\indent$y = (y_1, \dots, y_m) = S(x)$ — выходной вектор.

Используя скалярное произведение, это выражение можно записать в компактной форме:
\[
    (\alpha, x) \oplus \beta \cdot y = 0,
\]
где $(a,b) = a_1b_1 \oplus a_2b_2 \oplus \dots \oplus a_kb_k$ — скалярное произведение над $\mathbb{F}_2$.

\begin{definition}
Таблица линейных аппроксимаций (Linear Approximation Table, LAT) S-блока $S$ — это таблица размера $2^n \times 2^m$, элемент которой с индексами $[\alpha, \beta]$ определяется как:
\[
    \text{LAT}[\alpha, \beta] = |\{x \in \mathbb{F}_2^n \mid (\alpha, x) \oplus (\beta, S(x)) = 0\}| - 2^{n-1}
\]
\end{definition}

Значение $\text{LAT}[\alpha, \beta]$ показывает отклонение от равномерного распределения (смещение). Для идеального S-блока математическое ожидание $\text{LAT}[\alpha, \beta] = 0$.
\subsection{Критерий стойкости к линейному криптоанализу}

Эффективность линейной аппроксимации характеризуется её смещением. Чем больше абсолютное значение $|\text{LAT}[\alpha, \beta]|$, тем лучше аппроксимация для атакующего.

\begin{definition}
Линейность S-блока $S$ определяется как:
\[
    L(S) = \max_{\alpha, \beta \in \mathbb{F}_2^n \setminus \{0\}} |\text{LAT}[\alpha, \beta]|
\]
\end{definition}

\textbf{Критерий стойкости:} S-блок является стойким к линейному криптоанализу, если его линейность $L(S)$ минимальна.

\subsection{Аппроксимация как булева функция}

Рассмотрим линейную аппроксимацию как булеву функцию. Для фиксированных масок $\alpha$ и $\beta$ определим:
\[
    f_{\alpha,\beta}(x) =  (\alpha, x)  \oplus (\beta, S(x)) 
\]

Это булева функция $f_{\alpha,\beta} : \mathbb{F}_2^n \to \mathbb{F}_2$. 

Вес этой функции:
\[
    w(f_{\alpha,\beta}) = |\{x \in \mathbb{F}_2^n \mid f_{\alpha,\beta}(x) = 1\}|
\]

Связь с LAT:
\[
    \text{LAT}[\alpha, \beta] = 2^n - w(f_{\alpha, \beta}) - 2^{n-1} =2^{n-1} - w(f_{\alpha,\beta})
\]

\subsection{Нелинейность булевой функции}

\begin{definition}
Расстояние Хэмминга между функциями $f, g : \mathbb{F}_2^n \to \mathbb{F}_2$ определяется как:
\[
    d_H(f, g) = |\{x \in \mathbb{F}_2^n \mid f(x) \neq g(x)\}| = w(f \oplus g)
\]
\end{definition}

\begin{definition}
Нелинейность булевой функции $f : \mathbb{F}_2^n \to \mathbb{F}_2$ — это минимальное расстояние Хэмминга от $f$ до множества всех аффинных функций:
\[
    NL(f) = \min_{l \in \mathcal{A}_n} d(f, l) = 2^{n-1} - \frac{1}{2}\max_{a \in \mathbb{F}_2^n}|\widehat{f}(a)|,
\]
где $\mathcal{A}_n$ — множество всех аффинных функций от $n$ переменных; \\
\indent $\widehat{f}(\omega)$ — преобразование Уолша-Адамара функции $f$.
\end{definition}

Для S-блока $S : \mathbb{F}_2^n \to \mathbb{F}_2^m$ можно определить его покомпонентную нелинейность:
\[
    NL(S) = \min_{\beta \in \mathbb{F}_2^m \setminus \{0\}} NL(f_\beta),
\]
где $f_\beta(x) = (\beta, S(x))$ — булева функция, соответствующая линейной комбинации выходных координат.

\begin{theorem}
Связь между линейностью и нелинейностью S-блока:
\[
    L(S) = 2^{n-1} - NL(S)
\]
\end{theorem}

\begin{proof}
Рассмотрим булеву функцию $f_{\alpha, \beta}(x) = (\alpha, x) \oplus (\beta, S(x)) = (\alpha, x) \oplus f_\beta(x)$, где $f_\beta(x) = (\beta, S(x))$ — координатная функция S-блока.

\textbf{Шаг 1.} Если $f_{\alpha, \beta}(x) = 0$, то $f_\beta(x) = (\alpha, x)$. Если $f_{\alpha, \beta}(x) = 1$, то $f_\beta(x) = (\alpha, x) \oplus 1$.

\textbf{Шаг 2.} Расстояние Хэмминга от $f_\beta$ до аффинных функций $(\alpha, x)$ и $(\alpha, x) \oplus 1$:
\begin{align*}
    d_H(f_\beta, (\alpha, x)) &= wt(f_\beta \oplus (\alpha, x)) = wt(f_{\alpha, \beta}), \\
    d_H(f_\beta, (\alpha, x) \oplus 1) &= 2^n - d_H(f_\beta, (\alpha, x)) = 2^n - wt(f_{\alpha, \beta}).
\end{align*}

\textbf{Шаг 3.} Используя соотношение $w(f_{\alpha, \beta}) = 2^{n-1} - LAT[\alpha, \beta]$, получаем:
\begin{itemize}
    \item Если $LAT[\alpha, \beta] > 0$, то $wt(f_{\alpha, \beta}) < 2^{n-1}$, следовательно,
    \[
        \min\{d_H(f_\beta, (\alpha, x)), d_H(f_\beta, (\alpha, x) \oplus 1)\} = wt(f_{\alpha, \beta}) = 2^{n-1} - LAT[\alpha, \beta].
    \]
    
    \item Если $LAT[\alpha, \beta] < 0$, то $wt(f_{\alpha, \beta}) > 2^{n-1}$, следовательно,
    \[
        \min\{d_H(f_\beta, (\alpha, x)), d_H(f_\beta, (\alpha, x) \oplus 1)\} = 2^n - wt(f_{\alpha, \beta}) = 2^{n-1} + LAT[\alpha, \beta].
    \]
\end{itemize}

В обоих случаях минимальное расстояние равно $2^{n-1} - |LAT[\alpha, \beta]|$.

\textbf{Шаг 4.} Нелинейность функции $f_\beta$ — это минимальное расстояние до всех аффинных функций:
\[
    NL(f_\beta) = \min_{\substack{g \in \mathcal{A}_n}} d_H(f_\beta, g) = \min_{\alpha \in \mathbb{F}_2^n} \left(2^{n-1} - |LAT[\alpha, \beta]|\right) = 2^{n-1} - \max_{\alpha \neq 0} |LAT[\alpha, \beta]|,
\]
где $\mathcal{A}_n$ — множество всех аффинных функций от $n$ переменных.

\textbf{Шаг 5.} Нелинейность S-блока — это минимум по всем нетривиальным координатным функциям:
\begin{equation*}
\begin{aligned}
    NL(S) = \min_{\beta \neq 0} NL(f_\beta) = \min_{\beta \neq 0} \left(2^{n-1} - \max_{\alpha \neq 0} |LAT[\alpha, \beta]|\right) = \\
    2^{n-1} - \max_{\substack{\alpha \neq 0 \\ \beta \neq 0}} |LAT[\alpha, \beta]| = 2^{n-1} - L(S).
\end{aligned}
\end{equation*}
\end{proof}

\textbf{Следствие:} Максимизация нелинейности S-блока эквивалентна минимизации его линейности, что обеспечивает стойкость к линейному криптоанализу.


\subsection{Распределение нелинейностей координатных функций}

Для оценки качества S-блока рассматривают нелинейности всех нетривиальных линейных комбинаций выходов:
\[
    \{NL(f_\beta) \mid \beta \in \mathbb{F}_2^m \setminus \{0\}\}
\]

Помимо минимального значения $NL(S) = \min_{\beta \neq 0} NL(f_\beta)$, 
полезно вычислить:

\begin{itemize}
    \item \textbf{Среднюю нелинейность:}
    \[
        \overline{NL}(S) = \frac{1}{2^m - 1} \sum_{\beta \in \mathbb{F}_2^m \setminus \{0\}} NL(f_\beta)
    \]
    
    \item \textbf{Количество координатных функций с минимальной нелинейностью:}
    \[
        N_{\min} = |\{\beta \in \mathbb{F}_2^m \setminus \{0\} \mid NL(f_\beta) = NL(S)\}|
    \]
\end{itemize}

\textbf{Интерпретация:} 
\begin{itemize}
    \item Если $N_{\min}$ мало, то у атакующего меньше выбора хороших аппроксимаций
    \item Высокая средняя нелинейность $\overline{NL}(S)$ свидетельствует о равномерной стойкости
    \item Малая дисперсия нелинейностей указывает на сбалансированность S-блока
\end{itemize}
\newpage 
\section*{СПИСОК ИСПОЛЬЗОВАННЫХ ИСТОЧНИКОВ}
\addcontentsline{toc}{section}{СПИСОК ИСПОЛЬЗОВАННЫХ ИСТОЧНИКОВ}
\renewcommand{\refname}{\vspace*{-3em}}  % Уменьшаем пространство после заголовка

\begin{thebibliography}{9}
\renewcommand{\bibnumfmt}[1]{#1.}  % Убираем квадратные скобки из нумерации
\bibitem{biham1991}
Бихам, Э. Дифференциальный криптоанализ DES-подобных криптосистем [Текст] / Э. Бихам, А. Шамир. – Journal of Cryptology. – 1991. – Vol. 4, no. 1. – P. 3–72.

\bibitem{matsui1993}
Мацуи, М. Метод линейного криптоанализа для шифра DES [Текст] / М. Мацуи. – Advances in Cryptology – EUROCRYPT 1993. – LNCS 765. – Springer, 1994. – P. 386–397.

\bibitem{fips197}
NIST. FIPS 197: Advanced Encryption Standard (AES) [Текст]. – NIST, 2001. – 51 с.

\bibitem{gost2015}
ГОСТ Р 34.12-2015. Информационная технология. Криптографическая защита информации. Блочные шифры [Текст]. – М.: Стандартинформ, 2015.

\bibitem{biryukov2016}
Бирюков, А. Реверс-инжиниринг S-блока Стрибога, Кузнечика и STRIBOBr1 [Текст] / А. Бирюков, Л. Перрин, А. Удовенко. – Advances in Cryptology – EUROCRYPT 2016. – LNCS 9665. – Springer, 2016. – P. 372–402.

\bibitem{serpent1998}
Андерсон, Р. Serpent: предложение для стандарта AES [Текст] / Р. Андерсон, Э. Бихам, Л. Кнудсен. – AES Submission Document. – 1998. – 26 с.

\bibitem{twofish1998}
Шнайер, Б. Twofish: 128-битный блочный шифр [Текст] / Б. Шнайер и др. – AES Submission Document. – 1998. – 137 с.

\bibitem{sm4_2006}
OSCCA. GM/T 0002-2012. SM4 Block Cipher Algorithm [Текст]. – 2012.

\bibitem{camellia2001}
Аоки, К. Camellia: 128-битный блочный шифр, подходящий для разных платформ. Проектирование и анализ [Текст] / К. Аоки и др. – Selected Areas in Cryptography. – LNCS 2012. – Springer, 2001. – P. 39–56.

\bibitem{lightweight_cryptographic}
АНСАРИ, Ш. А., АЛИ, С. A systematic review of lightweight cryptographic schemes for security and privacy in IoT [Текст] // Discover Computing. – 2025. – Vol. 28, no. 1. – P. 1–57. – DOI: 10.1007/s10791-025-09755-3.
\end{thebibliography}




